\documentclass[10pt]{article}
\usepackage{graphicx} % Required for inserting images
\usepackage{url}
\usepackage{hyperref}
\title{IEG_Problems_Lecture9}
\author{martavictoriaperez }
\date{April 2026}

\usepackage[margin=1in]{geometry} 
\usepackage{amsmath,amsthm,amssymb, graphicx, multicol, array}
 
\newcommand{\N}{\mathbb{N}}
\newcommand{\Z}{\mathbb{Z}}
 
\newenvironment{problem}[2][Problem]{\begin{trivlist}
\item[\hskip \labelsep {\bfseries #1}\hskip \labelsep {\bfseries #2.}]}{\end{trivlist}}

\begin{document}
 
\title{\textbf{Lecture 9: Limiting CO$_2$ emissions}}
\author{
%Your name\\
DTU Course 46770: Integrated Energy Grids }
\maketitle

\begin{problem}{9.1}
Producing electricity in coal and gas power plants entails the efficiency, CO$_2$ emissions and fuel cost shown in Table \ref{Tab_coal_gas}. What is the cheapest option to produce electricity? Which CO$_2$ tax would be necessary to alter the merit order of the two technologies and incentivise lower CO$_2$ emissions per MWh of produced electricity?

\begin{table}[h]
    \centering
    \begin{tabular}{|c|c|c|}
        \hline
        & Coal Power plant	&Combined Cycle Gas Turbine (CCGT) \\ \hline
Fuel cost (USD/MWh thermal energy)	& 6.2	& 20.1 \\
Efficiency of power plant &	0.33	& 0.59  \\
(MWh electricity / MWh thermal energy) &	& \\
Emissions (tCO2/MWh thermal energy)	& 0.336	& 0.198 \\
\hline
    \end{tabular}
    \caption{Electricity production in coal and gas power plantns.}
    \label{Tab_coal_gas}
\end{table}
\end{problem}

\

\begin{problem}{9.2}
Build the model of Portugal described in Problem 3.1 in PyPSA and assume that methane gas emits 0.198 tCO2 per MWh of thermal energy contained in the gas. 

Consider the annualised capital costs and marginal generation costs for the different technologies in the following table.

\begin{table}[h]
    \centering
    \begin{tabular}{|c|c|c|}
    \hline
        Technology & Annualised capital costs (EUR/MW/a) & Marginal generation costs (EUR/MWh) \\
    \hline
    	Onshore Wind &  101,644 & 0 \\
         Solar PV &  51,346 & 0 \\
         OCGT & 47,718 &  64.7  \\
         CCGT & 104,788 &  46.8   \\
         Battery inverter &  24,678 & 0 \\
         Battery energy capacity &  12,894 & 0 \\
    \hline
    \end{tabular}
    \caption{Costs assumptions.}
    \label{tab:my_label}
\end{table}

Limit the maximum CO$_2$ emissions to 5 MtCO2/year. 

\begin{itemize}
\item[a)] Calculate the optimal installed capacities and plot the hourly generation and demand during January.
\item[b)] What is the CO$_2$ tax required to meet this CO$_2$ emission limit?

\end{itemize}

\end{problem}

\

\begin{problem}{9.3}

Build the model described in Problem 3.2 in PyPSA (i.e. the problem above but not including the CCGT generator with exogenous capacity).

\begin{itemize}
	\item[a)] Solve the problem for different CO$_2$ values ranging from 5 MtCO2/year to zero. Plot the total system cost and the required CO$_2$ prices as a function of the emissions allowance.
	\item[b)] Show that setting a CO$_2$ limit is equivalent to a carbon tax. Use your results from a) for the appropriate carbon taxes. Compare total system costs, emissions, and costs for carbon.
\end{itemize}



\end{problem}

%\begin{proof}[Solution]
%Write a solution here
%\end{proof}

\end{document}


 

